% Section 9.2, from lectures/L09/appendix/a_max_pool_layer.md, with the setup from the suite's
% README, lectures/L09/appendix/exercises/test/README.md. The solution to the code is in the
% repository.

\section{Exercises}
\label{c9:sec:exercises}\label{c9:app:a}

\begin{aside}
\textbf{How to work these.} The three sets are one class, written in the \code{cnn_work} codebase
you started in Chapter~\ref{c8:ch:conv}, \file{lectures/L08/cnn_work}. Keep working in the same
project: \file{max_pool.hpp} and \file{max_pool.cpp} are already in it as placeholders.

\textbf{How to check your work.} The test suite in \file{lectures/L09/appendix/exercises/test}
points back at \code{cnn_work}; Exercise~\ref{c9:ex:10} runs it, and \code{make ML_DIR=<path>}
points it at a tree of your own elsewhere.

The finished reference implementation is published in the companion repository as
\file{lectures/L10/cnn_demo}; see the warning below.
\end{aside}

\begin{warning}
\textbf{Spoiler warning.} \file{lectures/L10/cnn_demo} is the finished reference copy of this
project, and it contains working \code{Conv}, \code{MaxPool} and \code{Flatten} classes; i.e.\ the
answers to this chapter's exercises, the previous chapter's and the next. It's there for
Chapter~\ref{c10:ch:flatten}, where you compare your own pipeline against it. Write your own
version first; reading it now costs you the exercise.
\end{warning}

% ------------------------------------------------------------------------------------------
\exerciseset{Declaring the Max Pooling Layer}
\label{c9:set:1}
You'll continue extending \code{cnn_work}, implementing a concrete class \code{MaxPool}, replacing
the temporary stub class \code{ml::conv_layer::MaxPoolStub} currently used for the max pooling layer
in \file{factory.cpp}. This is the same logic you already worked through by hand in
\secref{c6:sec:exercises} and implemented as the standalone struct \code{ml::MaxPoolLayer} in
Chapter~\ref{c7:ch:layers}: here you're adapting that same logic to satisfy \code{cnn_work}'s
polymorphic \code{conv_layer::Interface} instead. \code{MaxPool} uses the \emph{same} interface as
\code{Conv} from Chapter~\ref{c8:ch:conv}; a max pooling layer just has no trainable parameters
and no activation function.

\codeexercise{The \code{MaxPool} class, declaration}
\label{c9:ex:1}
In the header file \file{include/ml/conv_layer/max_pool.hpp} (currently a placeholder), add a class
named \code{MaxPool} that inherits \code{ml::conv_layer::Interface}, the same interface
\code{Conv} implements, printed in Exercise~\ref{c8:ex:1}:
\begin{itemize}
  \item Use public inheritance and mark the class \code{final}.
  \item Override every method from the interface, including the destructor (which can be
    \code{= default}).
\end{itemize}

\codeexercise{Removed constructors and operators}
\label{c9:ex:2}
Delete the default constructor, copy constructor, move constructor, copy assignment operator, and
move assignment operator.

\codeexercise{Private member variables}
\label{c9:ex:3}
Add the following private member variables to \code{MaxPool}:
\begin{itemize}
  \item \textbf{\code{myInput}}: the most recently fed-forward input, needed later by
    \code{backpropagate()} to re-locate which position held the max value in each pool
    (\code{Matrix2d}).
  \item \textbf{\code{myInputGradients}}: input gradients; what \code{inputGradients()} returns
    (\code{Matrix2d}).
  \item \textbf{\code{myOutput}}: the pooled output; what \code{output()} returns
    (\code{Matrix2d}).
\end{itemize}
Note this differs from the standalone \code{ml::MaxPoolLayer} struct you wrote in
Chapter~\ref{c7:ch:layers}, which cached the max \emph{index} for each pool. Here we cache the
input and re-find the max during backpropagation, which keeps the member list smaller. Both give the
same answer, and both break ties toward the first occurrence in row-major order.

Feel free to add a private helper method that computes the pool size from \code{myInput} and
\code{myOutput} (\code{myInput.size() / myOutput.size()}) rather than storing it separately: that's
what the reference implementation does, since it's always derivable from the two matrices you
already have.

\codeexercise{Constructor}
\label{c9:ex:4}
Implement the constructor:

\begin{cppcode}
explicit MaxPool(std::size_t inputSize, std::size_t poolSize) noexcept;
\end{cppcode}

\begin{itemize}
  \item Print an error message and call \code{std::terminate()} if \code{inputSize} is \code{0},
    \code{poolSize} is \code{0}, \code{inputSize < poolSize}, or \code{inputSize} isn't evenly
    divisible by \code{poolSize}.
  \item Size \code{myInput} and \code{myInputGradients} to \code{inputSize}, and \code{myOutput} to
    \code{inputSize / poolSize} (see \code{ml::initMatrix()}).
\end{itemize}

\codeexercise{\code{feedforward()}}
\label{c9:ex:5}
Implement \code{bool feedforward(const Matrix2d& input) noexcept}:
\begin{itemize}
  \item Return \code{false} if the size of \code{input} doesn't match the size of \code{myInput}
    (\code{ml::matchDimensions()}), or if \code{input} isn't square (\code{ml::isMatrixSquare()}).
  \item For each output position \code{(i, j)}, find the largest value in the corresponding
    \code{poolSize × poolSize} block of \code{input} (the block starting at
    \code{(i * poolSize, j * poolSize)}) and store it in \code{myOutput[i][j]}. If a block has more
    than one occurrence of the max value, it doesn't matter which one \code{feedforward()} treats as
    ``the'' max; but \code{backpropagate()} (below) must be consistent with whichever one it picks,
    so route to the \emph{first} occurrence (which matches the convention of the hand-training
    exercise in \secref{c6:sec:exercises}).
  \item Store \code{input} in \code{myInput} (needed by \code{backpropagate()}).
  \item Return \code{true}.
\end{itemize}

\codeexercise{\code{backpropagate()}}
\label{c9:ex:6}
Implement \code{bool backpropagate(const Matrix2d& outputGradients) noexcept}:
\begin{itemize}
  \item Return \code{false} if the size of \code{outputGradients} doesn't match the size of
    \code{myOutput}, or if it isn't square.
  \item Reset \code{myInputGradients} to zero.
  \item For each output position \code{(i, j)}: re-locate which position in the corresponding block
    of \code{myInput} held the max value (first occurrence, same rule as \code{feedforward()}), and
    add \code{outputGradients[i][j]} to \code{myInputGradients} at that position. Every other
    position in the block gets no gradient (stays \code{0}).
  \item Return \code{true}.
\end{itemize}

\codeexercise{\code{optimize()}}
\label{c9:ex:7}
Implement \code{bool optimize(double learningRate) noexcept}. The layer has no trainable
parameters, so there's nothing to update: this is a genuine no-op, not a placeholder. Ignore
\code{learningRate} and return \code{true}.

\codeexercise{Trivial accessors}
\label{c9:ex:8}
Implement the four accessors from the interface:
\begin{itemize}
  \item \textbf{\code{inputSize()}}: returns \code{myInputGradients.size()}.
  \item \textbf{\code{outputSize()}}: returns \code{myOutput.size()}.
  \item \textbf{\code{output()}}: returns \code{myOutput}.
  \item \textbf{\code{inputGradients()}}: returns \code{myInputGradients}.
\end{itemize}

% ------------------------------------------------------------------------------------------
\exerciseset{Wiring It In}
\label{c9:set:2}

\codeexercise{Wiring it in}
\label{c9:ex:9}
\begin{enumerate}
  \item Make sure \file{source/ml/conv_layer/max_pool.cpp} is listed in \code{cnn_work}'s
    \file{Makefile} under \code{SRC_FILES}. It already is, as it has been since
    Chapter~\ref{c8:ch:conv}: the placeholder compiles to nothing.
  \item In \file{source/ml/factory/factory.cpp}, find \code{Factory::maxPoolLayer()}: it currently
    returns a \code{conv_layer::MaxPoolStub} (marked with a \code{//! @todo} comment). Replace that
    with your new class:
\begin{cppcode}
return std::make_unique<conv_layer::MaxPool>(inputSize, poolSize);
\end{cppcode}
  \item Run \code{make}. Both \code{Conv} and \code{MaxPool} are now real; only the flatten layer
    is still stubbed, so the output of \code{predict()} still won't be meaningful. That's the next
    chapter.
\end{enumerate}

% ------------------------------------------------------------------------------------------
\exerciseset{Running the Tests}
\label{c9:set:3}

\codeexercise{Running the tests}
\label{c9:ex:10}
The test suite is available in \file{lectures/L09/appendix/exercises/test}. It's cumulative: it
carries Chapter~\ref{c8:ch:conv}'s conv layer tests over unchanged and adds unit tests for
\code{MaxPool}.

Build and run it from this lecture's \file{appendix} directory:

\begin{shell}
make -C exercises/test
\end{shell}

\noindent There's nothing to copy: \code{ML_DIR} already points back at the \code{cnn_work}
codebase you started in Chapter~\ref{c8:ch:conv}. All 44 test cases should pass.

Because the pooling layer randomizes nothing, its expected values are written out in full rather
than recomputed the way the conv layer's are; and the main worked case is the hand-training example
from \secref{c6:sec:exercises}, the same pooling and gradient routing you did on paper.

The test worth knowing about before you start is \code{FeedforwardHandlesNegativeValues}. It feeds
an all-negative image, which catches starting the search for the max at \code{0.0} instead of at
the block's first value. Every non-negative input passes either way, so it's an easy bug to ship,
and it bites for real once the conv layer in front is using \code{Tanh}.

The test suite's README, \file{lectures/L09/appendix/exercises/test/README.md}, has more
information.
